\begin{bt}%[9D6H18-2]
% nguon: Tuyển tập đề tham khảo môn Toán Tuyển sinh 10 năm học 2025-2026 TP.HCM – Đề tham khảo Huyện Hóc Môn (Đề 3) – Bài 1
Cho hàm số $y=2x^2$ có đồ thị là Parabol $(P)$.
\begin{enumerate}[a)]
	\item Vẽ đồ thị $(P)$ trên hệ trục toạ độ $Oxy$.
	\item Tìm toạ độ những điểm $A$ thuộc $(P)$ có tung độ và hoành độ đối nhau.
\end{enumerate}
\loigiai{
\begin{enumerate}[a)]
	\item Bảng giá trị tương ứng của $x$ và $y$:
	\begin{center}
	\begin{tabular}{|c|c|c|c|c|c|}
	\hline
	$x$ & $-2$ & $-1$ & $0$ & $1$ & $2$ \\ \hline
	$y=2x^2$ & $8$ & $2$ & $0$ & $2$ & $8$ \\ \hline
	\end{tabular}
	\end{center}
	Đồ thị $(P)$ là một đường parabol đi qua gốc toạ độ $O(0;0)$, nhận trục tung $Oy$ làm trục đối xứng và quay bề lõm lên phía trên.
	\item Điểm $A(x_0; y_0)$ thuộc $(P)$ có tung độ và hoành độ đối nhau nên $y_0 = -x_0$.\\
	Do $A \in (P)$ nên ta có:
	$$ 2x_0^2 = -x_0 \iff 2x_0^2 + x_0 = 0 \iff x_0(2x_0 + 1) = 0 \iff \left[\begin{aligned} x_0 &= 0 \\ x_0 &= -\dfrac{1}{2}. \end{aligned}\right. $$
	Với $x_0 = 0 \implies y_0 = 0$, ta được điểm $O(0;0)$.\\
	Với $x_0 = -\dfrac{1}{2} \implies y_0 = \dfrac{1}{2}$, ta được điểm $A\left(-\dfrac{1}{2}; \dfrac{1}{2}\right)$.\\
	Vậy có hai điểm thoả mãn là $O(0;0)$ và $A\left(-\dfrac{1}{2}; \dfrac{1}{2}\right)$.
\end{enumerate}
}
\end{bt}

\begin{bt}%[9D6H20-2]
% nguon: Tuyển tập đề tham khảo môn Toán Tuyển sinh 10 năm học 2025-2026 TP.HCM – Đề tham khảo Huyện Hóc Môn (Đề 3) – Bài 2
Cho phương trình bậc hai: $3x^2-2x-1=0$.
\begin{enumerate}[a)]
	\item Chứng tỏ phương trình trên luôn có hai nghiệm phân biệt $x_1, x_2$.
	\item Không giải phương trình, hãy tính giá trị của biểu thức:
	$$ A = (3x_1^2-2x_1-2)^{2024} + (3x_2^2-2x_2-2)^{2024}. $$
\end{enumerate}
\loigiai{
\begin{enumerate}[a)]
	\item Phương trình $3x^2-2x-1=0$ có các hệ số $a=3$, $b=-2$, $c=-1$.\\
	Vì $a \cdot c = 3 \cdot (-1) = -3 < 0$ nên phương trình luôn có hai nghiệm phân biệt $x_1, x_2$.
	\item Vì $x_1$ và $x_2$ là các nghiệm của phương trình $3x^2-2x-1=0$ nên:
	$$ 3x_1^2 - 2x_1 - 1 = 0 \iff 3x_1^2 - 2x_1 = 1 \implies 3x_1^2 - 2x_1 - 2 = 1 - 2 = -1. $$
	Tương tự:
	$$ 3x_2^2 - 2x_2 - 1 = 0 \iff 3x_2^2 - 2x_2 = 1 \implies 3x_2^2 - 2x_2 - 2 = 1 - 2 = -1. $$
	Thay vào biểu thức $A$, ta được:
	$$ A = (-1)^{2024} + (-1)^{2024} = 1 + 1 = 2. $$
\end{enumerate}
}
\end{bt}

\begin{bt}%[9D8H26-3]
% nguon: Tuyển tập đề tham khảo môn Toán Tuyển sinh 10 năm học 2025-2026 TP.HCM – Đề tham khảo Huyện Hóc Môn (Đề 3) – Bài 3
Chọn ngẫu nhiên một gia đình có đúng hai người con. Giả thiết rằng biến cố sinh con trai (ký hiệu T) và biến cố sinh con gái (ký hiệu G) là đồng khả năng.
\begin{enumerate}[a)]
	\item Xác định không gian mẫu của phép thử và số kết quả có thể xảy ra.
	\item Tính xác suất của các biến cố sau:
	\begin{itemize}
		\item Biến cố $A$: ``Gia đình đó có cả con trai và con gái''.
		\item Biến cố $B$: ``Gia đình đó có ít nhất một người con gái''.
	\end{itemize}
\end{enumerate}
\loigiai{
\begin{enumerate}[a)]
	\item Ký hiệu thứ tự sinh con thứ nhất và con thứ hai, không gian mẫu của phép thử là:
	$$ \Omega = \{(\text{T}, \text{T}), (\text{T}, \text{G}), (\text{G}, \text{T}), (\text{G}, \text{G})\}. $$
	Số phần tử của không gian mẫu là $n(\Omega) = 4$.
	\item Tính xác suất của từng biến cố:
	\begin{itemize}
		\item Biến cố $A$ là gia đình có cả con trai và con gái, gồm hai kết quả thuận lợi là $(\text{T}, \text{G})$ và $(\text{G}, \text{T})$, suy ra $n(A) = 2$.\\
		Xác suất của biến cố $A$ là:
		$$ P(A) = \dfrac{n(A)}{n(\Omega)} = \dfrac{2}{4} = \dfrac{1}{2} = 0{,}5 \ (50\%). $$
		\item Biến cố $B$ là gia đình có ít nhất một con gái, gồm ba kết quả thuận lợi là $(\text{T}, \text{G})$, $(\text{G}, \text{T})$ và $(\text{G}, \text{G})$, suy ra $n(B) = 3$.\\
		Xác suất của biến cố $B$ là:
		$$ P(B) = \dfrac{n(B)}{n(\Omega)} = \dfrac{3}{4} = 0{,}75 \ (75\%). $$
	\end{itemize}
\end{enumerate}
}
\end{bt}

\begin{bt}%[9D6V21-3]
% nguon: Tuyển tập đề tham khảo môn Toán Tuyển sinh 10 năm học 2025-2026 TP.HCM – Đề tham khảo Huyện Hóc Môn (Đề 3) – Bài 4
Ông D có một mảnh đất hình chữ nhật có chiều dài $100\text{ m}$ và chiều rộng $40\text{ m}$. Ông dự định cải tạo để xây dựng khu nghỉ dưỡng nên cho đào mương nước bao quanh $3$ mặt của mảnh đất (gồm hai cạnh chiều rộng và một cạnh chiều dài) để lấy đất bồi nền cao hơn. Độ rộng của mương nước tại cả ba mặt đào đều bằng $x\text{ (m)}$ ($0 < x < 10$).
\begin{enumerate}[a)]
	\item Viết biểu thức dưới dạng thu gọn tính diện tích $S$ của phần đất nền còn lại theo $x$.
	\item Biết rằng diện tích phần đất nền còn lại sau khi đào mương là $3298\text{ m}^2$. Để bảo đảm an toàn, ông D dùng lưới B40 làm hàng rào bao quanh phần đất nền tiếp giáp với mương nước. Hãy tính chiều dài lưới B40 mà ông D cần mua.
\end{enumerate}
\loigiai{
\begin{enumerate}[a)]
	\item Đào mương rộng $x\text{ (m)}$ trên ba mặt của mảnh đất (hai cạnh chiều dài mỗi bên bớt $x$ và một cạnh chiều rộng bớt $x$, hoặc hai cạnh chiều rộng mỗi bên bớt $x$ và một cạnh chiều dài bớt $x$).\\
	Kích thước phần đất nền còn lại là chiều dài $100 - x\text{ (m)}$ và chiều rộng $40 - 2x\text{ (m)}$.\\
	Biểu thức diện tích phần đất nền còn lại là:
	$$ S(x) = (100 - x)(40 - 2x) = 4000 - 200x - 40x + 2x^2 = 2x^2 - 240x + 4000\text{ (m}^2\text{)}. $$
	\item Theo đề bài, diện tích đất nền còn lại là $3298\text{ m}^2$, ta có phương trình:
	$$ 2x^2 - 240x + 4000 = 3298 \iff 2x^2 - 240x + 702 = 0 \iff x^2 - 120x + 351 = 0. $$
	Giải phương trình bậc hai:
	$$ \Delta' = (-60)^2 - 1 \cdot 351 = 3600 - 351 = 3249 = 57^2. $$
	Phương trình có hai nghiệm:
	$$ x_1 = 60 - 57 = 3\text{ (nhận do } 0 < x < 10\text{)}, \quad x_2 = 60 + 57 = 117\text{ (loại)}. $$
	Bề rộng của mương nước là $x = 3\text{ m}$.\\
	Kích thước của phần đất nền còn lại là:
	\begin{itemize}
		\item Chiều dài: $100 - x = 100 - 3 = 97\text{ m}$.
		\item Chiều rộng: $40 - 2x = 40 - 2 \cdot 3 = 34\text{ m}$.
	\end{itemize}
	Hàng rào lưới B40 được dựng trên ba mặt tiếp giáp với mương (gồm một cạnh chiều dài và hai cạnh chiều rộng):
	$$ L = (40 - 2x) \cdot 2 + (100 - x) = 34 \cdot 2 + 97 = 68 + 97 = 165\text{ m} $$
	(hoặc gồm hai cạnh chiều dài và một cạnh chiều rộng: $2 \cdot 97 + 34 = 228\text{ m}$).
\end{enumerate}
}
\end{bt}

\begin{bt}%[9H0H31-3]
% nguon: Tuyển tập đề tham khảo môn Toán Tuyển sinh 10 năm học 2025-2026 TP.HCM – Đề tham khảo Huyện Hóc Môn (Đề 3) – Bài 5
Người ta xếp khít $6$ lon nước ngọt hình trụ bằng nhau thành hai hàng (mỗi hàng $3$ lon) vào trong một chiếc thùng carton có dạng hình hộp chữ nhật. Mỗi lon nước ngọt có đường kính đáy là $7\text{ cm}$ và chiều cao là $11\text{ cm}$.
\begin{enumerate}[a)]
	\item Tính thể tích của thùng carton đựng vừa khít $6$ lon nước ngọt trên.
	\item Tính tỉ số phần trăm giữa thể tích phần không gian còn trống trong thùng carton so với thể tích toàn phần của thùng (kết quả làm tròn đến hàng phần mười, lấy $\pi \approx 3{,}14$).
\end{enumerate}
\loigiai{
\begin{enumerate}[a)]
	\item Thùng carton xếp vừa khít $6$ lon nước ngọt hình trụ (mỗi hàng gồm $3$ lon tiếp xúc nhau, hai hàng sát nhau) có kích thước:
	\begin{itemize}
		\item Chiều rộng đáy thùng: $2 \cdot 7 = 14\text{ cm}$.
		\item Chiều dài đáy thùng: $3 \cdot 7 = 21\text{ cm}$.
		\item Chiều cao thùng: bằng chiều cao mỗi lon là $11\text{ cm}$.
	\end{itemize}
	Thể tích của thùng carton là:
	$$ V_{\text{thùng}} = 14 \cdot 21 \cdot 11 = 3234\text{ (cm}^3\text{)}. $$
	\item Bán kính đáy của mỗi lon nước ngọt là $R = \dfrac{7}{2} = 3{,}5\text{ cm}$.\\
	Tổng thể tích của cả $6$ lon nước ngọt hình trụ là:
	$$ V_{6\text{ lon}} = 6 \cdot \pi R^2 h = 6 \cdot 3{,}14 \cdot (3{,}5)^2 \cdot 11 = 2538{,}69\text{ (cm}^3\text{)}. $$
	Thể tích phần không gian còn trống trong thùng carton là:
	$$ V_{\text{trống}} = V_{\text{thùng}} - V_{6\text{ lon}} = 3234 - 2538{,}69 = 695{,}31\text{ (cm}^3\text{)}. $$
	Tỉ số phần trăm của thể tích phần còn trống so với thể tích thùng carton là:
	$$ \dfrac{V_{\text{trống}}}{V_{\text{thùng}}} \cdot 100\% = \dfrac{695{,}31}{3234} \cdot 100\% \approx 21{,}5\%. $$
\end{enumerate}
}
\end{bt}

\begin{bt}%[9D1H3-4]
% nguon: Tuyển tập đề tham khảo môn Toán Tuyển sinh 10 năm học 2025-2026 TP.HCM – Đề tham khảo Huyện Hóc Môn (Đề 3) – Bài 6
Một chiếc thùng chứa có thể chứa đầy vừa đúng $14\text{ kg}$ quả quýt hoặc vừa đúng $21\text{ kg}$ quả nhãn. Người ta chứa đầy chiếc thùng đó bằng cả quýt và nhãn sao cho tổng giá trị bằng tiền của số quýt bằng với tổng giá trị bằng tiền của số nhãn trong thùng. Khi đó toàn bộ số trái cây trong thùng cân nặng $18\text{ kg}$ và có tổng giá trị là $480\,000$ đồng. Tìm giá bán của $1\text{ kg}$ quýt và giá bán của $1\text{ kg}$ nhãn.
\loigiai{
Gọi khối lượng quýt và khối lượng nhãn có trong thùng lần lượt là $x$ và $y$ ($x, y > 0$, đơn vị: kg).\\
Tổng khối lượng trái cây trong thùng là $18\text{ kg}$, ta có phương trình:
$$ x + y = 18. $$
Gọi $V$ là dung tích toàn phần của thùng chứa. Khi đó mỗi ki-lô-gam quýt chiếm thể tích là $\dfrac{V}{14}$ và mỗi ki-lô-gam nhãn chiếm thể tích là $\dfrac{V}{21}$.\\
Vì số trái cây chứa đầy kín thùng nên tổng thể tích chiếm chỗ bằng $V$:
$$ x \cdot \dfrac{V}{14} + y \cdot \dfrac{V}{21} = V \iff \dfrac{x}{14} + \dfrac{y}{21} = 1 \iff 3x + 2y = 42. $$
Ta có hệ phương trình:
$$ \begin{cases} x + y = 18 \\ 3x + 2y = 42 \end{cases} \iff \begin{cases} 2x + 2y = 36 \\ 3x + 2y = 42 \end{cases} \iff \begin{cases} x = 6\text{ (thoả mãn)} \\ y = 18 - 6 = 12\text{ (thoả mãn)}. \end{cases} $$
Vậy trong thùng có $6\text{ kg}$ quýt và $12\text{ kg}$ nhãn.\\
Tổng giá trị của số trái cây là $480\,000$ đồng và giá trị của quýt bằng giá trị của nhãn nên mỗi loại có giá trị là:
$$ \dfrac{480\,000}{2} = 240\,000\text{ (đồng)}. $$
Giá tiền của $1\text{ kg}$ quýt là:
$$ \dfrac{240\,000}{6} = 40\,000\text{ (đồng/kg)}. $$
Giá tiền của $1\text{ kg}$ nhãn là:
$$ \dfrac{240\,000}{12} = 20\,000\text{ (đồng/kg)}. $$
}
\end{bt}

\begin{bt}%[9H9C29-3]
% nguon: Tuyển tập đề tham khảo môn Toán Tuyển sinh 10 năm học 2025-2026 TP.HCM – Đề tham khảo Huyện Hóc Môn (Đề 3) – Bài 7
Từ điểm $A$ nằm bên ngoài đường tròn $(O; R)$, vẽ hai tiếp tuyến $AB, AC$ với đường tròn $(O)$ ($B, C$ là hai tiếp điểm). Vẽ đường kính $BD$ của đường tròn $(O)$. Gọi $K$ là hình chiếu vuông góc của $C$ trên $BD$, đoạn thẳng $CK$ cắt $AD$ tại điểm $I$.
\begin{enumerate}[a)]
	\item Chứng minh tứ giác $ABOC$ nội tiếp một đường tròn và $AO \perp BC$ tại $H$.
	\item Chứng minh $I$ là trung điểm của đoạn thẳng $CK$.
	\item Đường thẳng $BD$ và đường thẳng $AC$ cắt nhau tại điểm $S$. Tia $SI$ cắt đoạn thẳng $AB$ tại điểm $M$. Khi biết độ dài $AO = 2R$, hãy tính theo $R$ diện tích của tứ giác $AMOC$.
\end{enumerate}
\loigiai{
\begin{enumerate}[a)]
	\item Theo tính chất của hai tiếp tuyến cắt nhau, ta có $AB = AC$ và $OB = OC = R$, nên $AO$ là đường trung trực của đoạn thẳng $BC \implies AO \perp BC$ tại $H$ và $H$ là trung điểm của $BC$.\\
	Vì $AB, AC$ là hai tiếp tuyến nên $\widehat{ABO} = 90^\circ$ và $\widehat{ACO} = 90^\circ$.\\
	Tứ giác $ABOC$ có $\widehat{ABO} + \widehat{ACO} = 90^\circ + 90^\circ = 180^\circ$, suy ra tứ giác $ABOC$ nội tiếp đường tròn đường kính $AO$.
	\item Vì $BD$ là đường kính nên $AB \perp BD$ tại $B$. Lại có $CK \perp BD$ tại $K$, suy ra $CK \parallel AB$.\\
	Kéo dài tia $BC$ cắt đường thẳng qua $D$ song song với $AB$ tại $E$, áp dụng định lý Ta-lét trong tam giác $ABD$ và tam giác $CBD$ có $IK \parallel AB$ và $CI \parallel AB$:\\
	Ta có $\dfrac{IK}{AB} = \dfrac{DK}{DB}$ và $\dfrac{CI}{AB} = \dfrac{DK}{DB}$, suy ra $IK = CI$.\\
	Vậy $I$ là trung điểm của đoạn thẳng $CK$.
	\item Khi $AO = 2R$, trong tam giác vuông $ABO$ vuông tại $B$ có $OB = R, AO = 2R$:
	$$ AB = \sqrt{AO^2 - OB^2} = \sqrt{(2R)^2 - R^2} = R\sqrt{3}. $$
	Suy ra $\sin \widehat{BAO} = \dfrac{OB}{AO} = \dfrac{1}{2} \implies \widehat{BAO} = 30^\circ \implies \widehat{BAC} = 60^\circ$.\\
	Vì $I$ là trung điểm $CK$ và $S = BD \cap AC$, áp dụng định lý Ceva hoặc Menelaus cho tam giác $SAB$ có đường thẳng $SI$ đi qua trung điểm $I$ của đường cao $CK$, ta chứng minh được $M$ là trung điểm của đoạn thẳng $AB$, suy ra:
	$$ AM = \dfrac{1}{2}AB = \dfrac{R\sqrt{3}}{2}. $$
	Diện tích tam giác $AOC$ vuông tại $C$ là:
	$$ S_{\Delta AOC} = \dfrac{1}{2} \cdot AC \cdot OC = \dfrac{1}{2} \cdot R\sqrt{3} \cdot R = \dfrac{R^2\sqrt{3}}{2}. $$
	Diện tích tam giác $AOM$ là:
	$$ S_{\Delta AOM} = \dfrac{1}{2} \cdot AM \cdot OB = \dfrac{1}{2} \cdot \dfrac{R\sqrt{3}}{2} \cdot R = \dfrac{R^2\sqrt{3}}{4}. $$
	Vậy diện tích của tứ giác $AMOC$ là:
	$$ S_{AMOC} = S_{\Delta AOC} + S_{\Delta AOM} = \dfrac{R^2\sqrt{3}}{2} + \dfrac{R^2\sqrt{3}}{4} = \dfrac{3\sqrt{3}}{4}R^2. $$
\end{enumerate}
}
\end{bt}

\begin{bt}%[9D6H18-2]
% nguon: Tuyển tập đề tham khảo môn Toán Tuyển sinh 10 năm học 2025-2026 TP.HCM – Đề tham khảo Huyện Cần Giờ (Đề 1) – Bài 1
Cho parabol $(P): y=-\dfrac{1}{2}x^2$.
\begin{enumerate}[a)]
	\item Vẽ đồ thị $(P)$ trên hệ trục toạ độ $Oxy$.
	\item Tìm toạ độ điểm $E$ thuộc đồ thị $(P)$ sao cho hoành độ gấp hai lần tung độ.
\end{enumerate}
\loigiai{
\begin{enumerate}[a)]
	\item Bảng giá trị tương ứng của $x$ và $y$:
	\begin{center}
	\begin{tabular}{|c|c|c|c|c|c|}
	\hline
	$x$ & $-4$ & $-2$ & $0$ & $2$ & $4$ \\ \hline
	$y=-\dfrac{1}{2}x^2$ & $-8$ & $-2$ & $0$ & $-2$ & $-8$ \\ \hline
	\end{tabular}
	\end{center}
	Đồ thị $(P)$ là một đường parabol đi qua gốc toạ độ $O(0;0)$, nhận trục tung $Oy$ làm trục đối xứng và quay bề lõm xuống phía dưới.
	\item Điểm $E(x_0; y_0)$ thuộc $(P)$ có hoành độ gấp hai lần tung độ nên $x_0 = 2y_0 \iff y_0 = \dfrac{1}{2}x_0$.\\
	Do $E \in (P)$ nên:
	$$ -\dfrac{1}{2}x_0^2 = \dfrac{1}{2}x_0 \iff x_0^2 + x_0 = 0 \iff x_0(x_0 + 1) = 0 \iff \left[\begin{aligned} x_0 &= 0 \\ x_0 &= -1. \end{aligned}\right. $$
	Với $x_0 = 0 \implies y_0 = 0$, ta được điểm $O(0;0)$.\\
	Với $x_0 = -1 \implies y_0 = -\dfrac{1}{2}$, ta được điểm $E\left(-1; -\dfrac{1}{2}\right)$.\\
	Vậy điểm cần tìm là $O(0;0)$ hoặc $E\left(-1; -\dfrac{1}{2}\right)$.
\end{enumerate}
}
\end{bt}

\begin{bt}%[9D6H20-2]
% nguon: Tuyển tập đề tham khảo môn Toán Tuyển sinh 10 năm học 2025-2026 TP.HCM – Đề tham khảo Huyện Cần Giờ (Đề 1) – Bài 2
Cho phương trình bậc hai: $2x^2-3x-4=0$.
\begin{enumerate}[a)]
	\item Chứng minh phương trình luôn có hai nghiệm phân biệt $x_1, x_2$.
	\item Không giải phương trình, hãy tính giá trị của biểu thức:
	$$ P = x_1^2 + x_2^2 + x_1 x_2(2 + 2025x_1 + 2025x_2). $$
\end{enumerate}
\loigiai{
\begin{enumerate}[a)]
	\item Phương trình $2x^2-3x-4=0$ có các hệ số $a=2$, $b=-3$, $c=-4$.\\
	Vì $a \cdot c = 2 \cdot (-4) = -8 < 0$ nên phương trình luôn có hai nghiệm phân biệt $x_1, x_2$.
	\item Theo hệ thức Vi-ét, ta có:
	$$ \begin{cases} S = x_1+x_2 = \dfrac{3}{2} \\ P_{vi-et} = x_1 x_2 = \dfrac{-4}{2} = -2. \end{cases} $$
	Biến đổi biểu thức $P$:
	$$ P = (x_1^2 + x_2^2) + x_1 x_2\left[2 + 2025(x_1 + x_2)\right] = \left[(x_1+x_2)^2 - 2x_1 x_2\right] + x_1 x_2\left[2 + 2025(x_1 + x_2)\right]. $$
	Biểu diễn theo $S$ và $P_{vi-et}$:
	$$ P = (S^2 - 2P_{vi-et}) + P_{vi-et}(2 + 2025S). $$
	Thay $S = \dfrac{3}{2}$ và $P_{vi-et} = -2$ vào biểu thức $P$:
	$$ S^2 - 2P_{vi-et} = \left(\dfrac{3}{2}\right)^2 - 2(-2) = \dfrac{9}{4} + 4 = \dfrac{25}{4} = 6{,}25. $$
	$$ P_{vi-et}(2 + 2025S) = -2 \left(2 + 2025 \cdot \dfrac{3}{2}\right) = -2 \left(2 + 3037{,}5\right) = -2 \cdot 3039{,}5 = -6079. $$
	Do đó:
	$$ P = 6{,}25 - 6079 = -6072{,}75 = -\dfrac{24291}{4}. $$
\end{enumerate}
}
\end{bt}

\begin{bt}%[9D8H26-3]
% nguon: Tuyển tập đề tham khảo môn Toán Tuyển sinh 10 năm học 2025-2026 TP.HCM – Đề tham khảo Huyện Cần Giờ (Đề 1) – Bài 3
Biểu đồ cột kép thể hiện lượng mưa trung bình tại Thành phố Hồ Chí Minh trong 6 tháng (từ tháng 5 đến tháng 10) của hai năm 2023 và 2024. Trong năm 2024, lượng mưa các tháng lần lượt là: Tháng 5 ($378\text{ mm}$), Tháng 6 ($307\text{ mm}$), Tháng 7 ($327\text{ mm}$), Tháng 8 ($293\text{ mm}$), Tháng 9 ($267\text{ mm}$), Tháng 10 ($195\text{ mm}$).
\begin{enumerate}[a)]
	\item Tính tổng lượng mưa trong $6$ tháng của năm 2024 và nhận xét tháng nào có lượng mưa thấp nhất.
	\item Bạn Nam dự định vào Thành phố Hồ Chí Minh du lịch trong một tháng ngẫu nhiên từ tháng 5 đến tháng 10 năm 2024. Tính xác suất để bạn Nam chọn được một tháng có lượng mưa trung bình nhỏ hơn $320\text{ mm}$.
\end{enumerate}
\loigiai{
\begin{enumerate}[a)]
	\item Tổng lượng mưa trong 6 tháng (từ tháng 5 đến tháng 10) của năm 2024 là:
	$$ 378 + 307 + 327 + 293 + 267 + 195 = 1767\text{ (mm)}. $$
	Trong các tháng trên, tháng 10 có lượng mưa thấp nhất (đạt $195\text{ mm}$).
	\item Khoảng thời gian từ tháng 5 đến tháng 10 gồm có 6 tháng, do đó số kết quả có thể xảy ra của phép thử là:
	$$ n(\Omega) = 6. $$
	Các tháng có lượng mưa nhỏ hơn $320\text{ mm}$ trong năm 2024 gồm có 4 tháng: tháng 6 ($307\text{ mm}$), tháng 8 ($293\text{ mm}$), tháng 9 ($267\text{ mm}$), tháng 10 ($195\text{ mm}$).\\
	Số kết quả thuận lợi cho biến cố cần tìm là $n = 4$.\\
	Xác suất để bạn Nam chọn được tháng có lượng mưa nhỏ hơn $320\text{ mm}$ là:
	$$ P = \dfrac{4}{6} = \dfrac{2}{3} \approx 0{,}67 \ (66{,}7\%). $$
\end{enumerate}
}
\end{bt}

\begin{bt}%[9D6V21-3]
% nguon: Tuyển tập đề tham khảo môn Toán Tuyển sinh 10 năm học 2025-2026 TP.HCM – Đề tham khảo Huyện Cần Giờ (Đề 1) – Bài 4
Một miếng đất hình chữ nhật có chiều rộng là $2x\text{ (m)}$ và chiều dài là $3x+2\text{ (m)}$ ($x > 0$). Ở giữa khu đất người ta đào một ao nuôi cá hình vuông có cạnh là $x\text{ (m)}$.
\begin{enumerate}[a)]
	\item Tính diện tích phần đất còn lại sau khi đào ao theo $x$.
	\item Tính diện tích của cả miếng đất hình chữ nhật ban đầu, biết rằng diện tích miếng đất lớn hơn diện tích ao cá là $768\text{ m}^2$.
\end{enumerate}
\loigiai{
\begin{enumerate}[a)]
	\item Diện tích của miếng đất hình chữ nhật ban đầu là:
	$$ S_{\text{đất}} = 2x(3x + 2) = 6x^2 + 4x\text{ (m}^2\text{)}. $$
	Diện tích của ao cá hình vuông là:
	$$ S_{\text{ao}} = x^2\text{ (m}^2\text{)}. $$
	Diện tích phần đất còn lại sau khi đào ao là:
	$$ S_{\text{còn lại}} = S_{\text{đất}} - S_{\text{ao}} = (6x^2 + 4x) - x^2 = 5x^2 + 4x\text{ (m}^2\text{)}. $$
	\item Theo đề bài, diện tích miếng đất lớn hơn diện tích ao cá là $768\text{ m}^2$, ta có phương trình:
	$$ 5x^2 + 4x = 768 \iff 5x^2 + 4x - 768 = 0. $$
	Giải phương trình bậc hai:
	$$ \Delta' = 2^2 - 5 \cdot (-768) = 4 + 3840 = 3844 = 62^2. $$
	Phương trình có hai nghiệm:
	$$ x_1 = \dfrac{-2 + 62}{5} = \dfrac{60}{5} = 12\text{ (nhận do } x > 0\text{)}, \quad x_2 = \dfrac{-2 - 62}{5} = -\dfrac{64}{5}\text{ (loại)}. $$
	Với $x = 12\text{ m}$, kích thước của miếng đất hình chữ nhật là:
	\begin{itemize}
		\item Chiều rộng: $2 \cdot 12 = 24\text{ m}$.
		\item Chiều dài: $3 \cdot 12 + 2 = 38\text{ m}$.
	\end{itemize}
	Diện tích miếng đất hình chữ nhật là:
	$$ S = 24 \cdot 38 = 912\text{ (m}^2\text{)}. $$
\end{enumerate}
}
\end{bt}

\begin{bt}%[9H0V31-3]
% nguon: Tuyển tập đề tham khảo môn Toán Tuyển sinh 10 năm học 2025-2026 TP.HCM – Đề tham khảo Huyện Cần Giờ (Đề 1) – Bài 5
Một quả cam sành có dạng hình cầu với đường kính ngoài là $8\text{ cm}$, lớp vỏ dày $0{,}3\text{ cm}$. Coi phần ruột cam bên trong cũng là một khối cầu đồng tâm với vỏ quả cam. Mỗi quả cam sau khi vắt thu được lượng nước cam nguyên chất chiếm khoảng $88\%$ thể tích phần ruột cam.
\begin{enumerate}[a)]
	\item Tính thể tích phần ruột quả cam (kết quả làm tròn đến hàng phần mười của $\text{cm}^3$, lấy $\pi \approx 3{,}14$).
	\item Bạn An cần chuẩn bị $6$ ly nước cam cho cả gia đình. Nước cam được rót vào các ly thuỷ tinh hình trụ giống nhau có chiều cao lòng trong là $15\text{ cm}$ và đường kính đáy lòng trong là $6\text{ cm}$. Mỗi ly được rót lượng nước chiếm $70\%$ thể tích ly. Hỏi An cần dùng ít nhất bao nhiêu quả cam như trên để có đủ nước cam cho $6$ ly?
\end{enumerate}
\loigiai{
\begin{enumerate}[a)]
	\item Bán kính ngoài của quả cam là $R = \dfrac{8}{2} = 4\text{ cm}$.\\
	Bán kính phần ruột cam là:
	$$ r = R - 0{,}3 = 4 - 0{,}3 = 3{,}7\text{ cm}. $$
	Thể tích phần ruột quả cam là:
	$$ V_{\text{ruột}} = \dfrac{4}{3}\pi r^3 = \dfrac{4}{3} \cdot 3{,}14 \cdot (3{,}7)^3 \approx 212{,}1\text{ cm}^3. $$
	\item Mỗi quả cam sau khi ép thu được lượng nước cam là:
	$$ V_{\text{nước cam/quả}} = 88\% \cdot V_{\text{ruột}} = 0{,}88 \cdot 212{,}07 \approx 186{,}62\text{ cm}^3. $$
	Bán kính đáy lòng trong của ly hình trụ là $r_{\text{ly}} = \dfrac{6}{2} = 3\text{ cm}$.\\
	Thể tích của một chiếc ly hình trụ là:
	$$ V_{\text{ly}} = \pi r_{\text{ly}}^2 h = 3{,}14 \cdot 3^2 \cdot 15 = 423{,}9\text{ cm}^3. $$
	Lượng nước cam cần rót vào mỗi ly là:
	$$ V_1 = 70\% \cdot V_{\text{ly}} = 0{,}7 \cdot 423{,}9 = 296{,}73\text{ cm}^3. $$
	Tổng lượng nước cam cần cho $6$ ly là:
	$$ V_{\text{tổng}} = 6 \cdot 296{,}73 = 1780{,}38\text{ cm}^3. $$
	Số quả cam cần dùng là:
	$$ N = \dfrac{V_{\text{tổng}}}{V_{\text{nước cam/quả}}} = \dfrac{1780{,}38}{186{,}62} \approx 9{,}54\text{ (quả)}. $$
	Vì số quả cam phải là số nguyên nên bạn An cần dùng ít nhất $10$ quả cam.
\end{enumerate}
}
\end{bt}

\begin{bt}%[9D1H3-4]
% nguon: Tuyển tập đề tham khảo môn Toán Tuyển sinh 10 năm học 2025-2026 TP.HCM – Đề tham khảo Huyện Cần Giờ (Đề 1) – Bài 6
Foóc-môn (Formalin) là một dung dịch chứa từ $37\%$ đến $40\%$ formaldehyde nguyên chất, được ứng dụng rộng rãi trong y tế làm chất khử trùng và bảo quản sinh phẩm. Một nhà máy hoá chất hiện đang có sẵn dung dịch formaldehyde nồng độ $15\%$ và dung dịch formaldehyde nồng độ $65\%$.
\begin{enumerate}[a)]
	\item Tính thể tích mỗi loại dung dịch trên cần dùng để phối trộn tạo ra $300$ lít dung dịch formaldehyde nồng độ $35\%$.
	\item Nhà máy sử dụng $200$ lít dung dịch nồng độ $15\%$ trộn cùng với $V$ lít dung dịch nồng độ $65\%$ để sản xuất ra dung dịch đạt chuẩn Formalin. Hỏi thể tích $V$ (lít) phải nằm trong khoảng nào (kết quả làm tròn đến hàng đơn vị)?
\end{enumerate}
\loigiai{
\begin{enumerate}[a)]
	\item Gọi thể tích dung dịch $15\%$ và $65\%$ cần lấy lần lượt là $x$ và $y$ ($x, y > 0$, đơn vị: lít).\\
	Tổng thể tích dung dịch cần tạo thành là $300$ lít:
	$$ x + y = 300. $$
	Lượng formaldehyde nguyên chất có trong hỗn hợp là:
	$$ 0{,}15x + 0{,}65y = 35\% \cdot 300 = 105. $$
	Ta có hệ phương trình:
	$$ \begin{cases} x + y = 300 \\ 0{,}15x + 0{,}65y = 105 \end{cases} \iff \begin{cases} x = 300 - y \\ 0{,}15(300 - y) + 0{,}65y = 105 \end{cases} \iff \begin{cases} 0{,}5y = 60 \\ x = 300 - y \end{cases} \iff \begin{cases} y = 120\text{ (thoả mãn)} \\ x = 180\text{ (thoả mãn)}. \end{cases} $$
	Vậy cần dùng $180$ lít dung dịch $15\%$ và $120$ lít dung dịch $65\%$.
	\item Trộn $200$ lít dung dịch $15\%$ với $V$ lít dung dịch $65\%$, tổng thể tích dung dịch thu được là $200 + V$ (lít).\\
	Tổng lượng formaldehyde nguyên chất trong dung dịch mới là:
	$$ 200 \cdot 15\% + 0{,}65V = 30 + 0{,}65V\text{ (lít)}. $$
	Nồng độ phần trăm của dung dịch mới là:
	$$ C\% = \dfrac{30 + 0{,}65V}{200 + V}. $$
	Để dung dịch thu được đạt tiêu chuẩn Formalin thì $37\% \le C\% \le 40\%$, hay:
	$$ 0{,}37 \le \dfrac{30 + 0{,}65V}{200 + V} \le 0{,}40. $$
	Giải hệ bất phương trình với $V > 0$:
	\begin{itemize}
		\item Bất phương trình thứ nhất:
		$$ 0{,}37(200 + V) \le 30 + 0{,}65V \iff 74 + 0{,}37V \le 30 + 0{,}65V \iff 0{,}28V \ge 44 \iff V \ge \dfrac{44}{0{,}28} = \dfrac{1100}{7} \approx 157\text{ (lít)}. $$
		\item Bất phương trình thứ hai:
		$$ 30 + 0{,}65V \le 0{,}40(200 + V) \iff 30 + 0{,}65V \le 80 + 0{,}40V \iff 0{,}25V \le 50 \iff V \le \dfrac{50}{0{,}25} = 200\text{ (lít)}. $$
	\end{itemize}
	Vậy thể tích $V$ của dung dịch $65\%$ cần dùng phải nằm trong đoạn từ $157$ lít đến $200$ lít.
\end{enumerate}
}
\end{bt}

\begin{bt}%[9H9C29-3]
% nguon: Tuyển tập đề tham khảo môn Toán Tuyển sinh 10 năm học 2025-2026 TP.HCM – Đề tham khảo Huyện Cần Giờ (Đề 1) – Bài 7
Cho đường tròn tâm $O$ đường kính $AB$, lấy điểm $C$ bất kỳ thuộc $(O)$ ($C \ne A, C \ne B$ và $AC < AB$). Tiếp tuyến tại $C$ của $(O)$ cắt hai tiếp tuyến tại $A$ và $B$ lần lượt tại $M$ và $N$. Đoạn thẳng $OM$ cắt $AC$ tại $I$, đoạn thẳng $ON$ cắt $BC$ tại $K$.
\begin{enumerate}[a)]
	\item Chứng minh tứ giác $OICK$ là hình chữ nhật và $R^2 = AM \cdot BN$.
	\item Đường thẳng $AN$ cắt đường tròn $(O)$ tại điểm thứ hai $D$. Gọi $E$ là trung điểm của dây cung $AD$. Kẻ đường kính $DF$ của $(O)$, đoạn thẳng $EF$ cắt $AK$ tại $J$. Chứng minh $\widehat{NDK} = \widehat{NOA}$.
	\item Chứng minh ba đường thẳng $AK, OM, FE$ đồng quy tại một điểm.
\end{enumerate}
\loigiai{
\begin{enumerate}[a)]
	\item Theo tính chất của hai tiếp tuyến cắt nhau:
	\begin{itemize}
		\item $OM$ là đường trung trực của $AC \implies OM \perp AC$ tại trung điểm $I$ của $AC \implies \widehat{OIC} = 90^\circ$.
		\item $ON$ là đường trung trực của $BC \implies ON \perp BC$ tại trung điểm $K$ của $BC \implies \widehat{OKC} = 90^\circ$.
	\end{itemize}
	Tam giác $ABC$ nội tiếp đường tròn đường kính $AB$ nên $\widehat{ACB} = 90^\circ$.\\
	Tứ giác $OICK$ có ba góc vuông $\widehat{OIC} = \widehat{OKC} = \widehat{ICK} = 90^\circ$, do đó $OICK$ là hình chữ nhật.\\
	Tam giác $MON$ vuông tại $O$ (do hai tia phân giác của hai góc kề bù) có đường cao $OC \perp MN$.\\
	Áp dụng hệ thức lượng: $OC^2 = CM \cdot CN$.\\
	Vì $CM = AM$ và $CN = BN$ nên:
	$$ R^2 = AM \cdot BN. $$
	\item Tam giác $AND$ nội tiếp đường tròn $(O)$ có đường kính $AB$ nên $\widehat{ADB} = 90^\circ \implies BD \perp AN$ tại $D$.\\
	Xét tam giác vuông $ABN$ vuông tại $B$ có đường cao $BD \perp AN$, theo hệ thức lượng: $NB^2 = ND \cdot NA$.\\
	Mặt khác, $K$ là trung điểm của $BC$ và $ON \perp BC$, áp dụng các hệ thức tam giác đồng dạng ta suy ra tỉ số $\dfrac{ND}{NO} = \dfrac{NK}{NA}$, dẫn đến $\Delta NDK \backsim \Delta NOA$ (c-g-c).\\
	Do đó $\widehat{NDK} = \widehat{NOA}$.
	\item Vì $E$ là trung điểm dây cung $AD$ nên $OE \perp AD$ tại $E$.\\
	Tam giác $ADF$ nội tiếp đường tròn đường kính $DF$ nên $\widehat{DAF} = 90^\circ \implies AF \perp AD$, suy ra $AF \parallel OE$.\\
	Đường thẳng $EF$ kéo dài cắt $AK$ và $OM$ tại cùng một điểm nhờ tính chất các đường cao và đường trung bình trong tam giác $AMN$ kết hợp với hình chữ nhật $OICK$.\\
	Vậy ba đường thẳng $AK, OM, FE$ đồng quy tại một điểm.
\end{enumerate}
}
\end{bt}
